\section{Probabilités}
\label{sec:probabilites}

Les données sont vues comme les réalisations de variables aléatoires. C'est ce qui permet de parler d'« erreur moyenne » d'un modèle sur toute une population (chapitre~1).

\subsection{Variables aléatoires, espérance, variance}

\begin{definition}[Variable aléatoire, loi]
Une \emph{variable aléatoire} $X$ est une quantité dont la valeur dépend du hasard (l'individu tiré au sort, par exemple).
Elle est \emph{discrète} si elle prend des valeurs isolées $x_1,x_2,\dots$ avec des probabilités $p_k=\P(X=x_k)$ ;
elle est \emph{continue} si elle a une \emph{densité} $f\ge0$ telle que $\P(a\le X\le b)=\int_a^bf(x)\,dx$.
\end{definition}

\begin{definition}[Espérance, variance, écart-type]
L'\emph{espérance} (moyenne théorique) de $X$ est $\E[X]=\sum_kx_kp_k$ dans le cas discret, $\E[X]=\int xf(x)\,dx$ dans le cas continu.
La \emph{variance} mesure la dispersion autour de la moyenne : $\Var(X)=\E\big[(X-\E[X])^2\big]$ ; l'\emph{écart-type} est $\sqrt{\Var(X)}$.
\end{definition}

\begin{proposition}[Règles de calcul]
\label{prop:regles}
Pour des variables $X,Y$ et des réels $a,b$ :
\begin{enumerate}[label=(\roman*)]
  \item $\E[aX+bY]=a\E[X]+b\E[Y]$ (linéarité) ;
  \item $\Var(X)=\E[X^2]-(\E[X])^2$ et $\Var(aX+b)=a^2\Var(X)$.
\end{enumerate}
\end{proposition}
\begin{preuve}
(i) est admis. (ii) : en développant, $\E[(X-\E X)^2]=\E[X^2]-2\E[X]\E[X]+(\E X)^2$. Puis $aX+b-\E[aX+b]=a(X-\E X)$, dont le carré vaut $a^2(X-\E X)^2$.
\end{preuve}

\begin{proposition}[La meilleure constante est l'espérance]
\label{prop:meilleureconstante}
Pour toute variable $Z$ de variance finie et tout réel $c$ :
\[
  \E\big[(Z-c)^2\big]=\Var(Z)+\big(\E[Z]-c\big)^2 .
\]
Donc $c\mapsto\E[(Z-c)^2]$ atteint son minimum en $c=\E[Z]$, et la valeur minimale est $\Var(Z)$.
\end{proposition}
\begin{preuve}
On écrit $Z-c=(Z-\E Z)+(\E Z-c)$ et l'on développe :
$\E[(Z-c)^2]=\Var(Z)+2(\E Z-c)\E[Z-\E Z]+(\E Z-c)^2$, et $\E[Z-\E Z]=0$.
\end{preuve}
C'est la version « population » de la proposition~\ref{prop:moyenne}, et c'est elle qui donne le meilleur prédicteur au chapitre~1.

\subsection{Covariance, corrélation, indépendance}

\begin{definition}[Covariance, corrélation]
$\Cov(X,Y)=\E\big[(X-\E X)(Y-\E Y)\big]$ et, si les variances sont non nulles, la \emph{corrélation} est
$\rho(X,Y)=\dfrac{\Cov(X,Y)}{\sqrt{\Var(X)\Var(Y)}}$, toujours comprise entre $-1$ et $1$ (admis).
\end{definition}

\begin{definition}[Indépendance]
$X$ et $Y$ sont \emph{indépendantes} si connaître la valeur de l'une ne change pas la loi de l'autre : $\P(X\in A,\,Y\in B)=\P(X\in A)\,\P(Y\in B)$ pour tous ensembles $A$, $B$.
\end{definition}

\begin{proposition}[admise]
\label{prop:indep}
Si $X$ et $Y$ sont indépendantes, alors $\Cov(X,Y)=0$. Dans tous les cas, $\Var(X+Y)=\Var(X)+\Var(Y)+2\Cov(X,Y)$.
\end{proposition}

\begin{remarque}
Une covariance nulle n'implique pas l'indépendance : la corrélation ne mesure que la liaison \emph{linéaire}.
\end{remarque}

\subsection{Lois usuelles}

\begin{center}
\renewcommand{\arraystretch}{1.3}
\begin{tabular}{llll}
\toprule
\textbf{Loi} & \textbf{Description} & $\E[X]$ & $\Var(X)$ \\
\midrule
Bernoulli $\mathcal B(p)$ & $X=1$ avec probabilité $p$, $0$ sinon & $p$ & $p(1-p)$ \\
Binomiale $\mathcal B(n,p)$ & nombre de succès en $n$ essais indépendants & $np$ & $np(1-p)$ \\
Uniforme $\mathcal U[a,b]$ & densité constante $\frac{1}{b-a}$ sur $[a,b]$ & $\frac{a+b}{2}$ & $\frac{(b-a)^2}{12}$ \\
Normale $\mathcal N(\mu,\sigma^2)$ & densité $\frac{1}{\sigma\sqrt{2\pi}}e^{-(x-\mu)^2/(2\sigma^2)}$ & $\mu$ & $\sigma^2$ \\
\bottomrule
\end{tabular}
\end{center}

Si $Z\sim\mathcal N(0,1)$, on note $\Phi(t)=\P(Z\le t)$ ; on retiendra $\P(-1{,}96\le Z\le1{,}96)\approx0{,}95$.

\subsection{Moyenne d'un échantillon et théorèmes limites}

\begin{proposition}[Espérance et variance d'une moyenne]
\label{prop:moyennevar}
Soient $X_1,\dots,X_n$ indépendantes et de même loi, de moyenne $\mu$ et de variance $\sigma^2$, et $\bar X_n=\frac1n\sum_iX_i$. Alors
$\E[\bar X_n]=\mu$ et $\Var(\bar X_n)=\dfrac{\sigma^2}{n}$.
\end{proposition}
\begin{preuve}
Par linéarité, $\E[\bar X_n]=\frac1n\cdot n\mu=\mu$. Par indépendance (proposition~\ref{prop:indep}), les covariances sont nulles et
$\Var(\bar X_n)=\frac{1}{n^2}\sum_i\Var(X_i)=\frac{n\sigma^2}{n^2}=\frac{\sigma^2}{n}$.
\end{preuve}

\begin{theoreme}[Loi des grands nombres, admis]
\label{thm:lgn}
Sous les mêmes hypothèses, $\bar X_n$ se rapproche de $\mu$ quand $n$ devient grand : pour tout $\varepsilon>0$, $\P(|\bar X_n-\mu|>\varepsilon)\to0$.
\end{theoreme}

\begin{theoreme}[Théorème central limite, admis]
\label{thm:tcl}
Pour $n$ grand, $\bar X_n$ suit approximativement la loi $\mathcal N\big(\mu,\sigma^2/n\big)$. En particulier,
$\P\big(|\bar X_n-\mu|\le1{,}96\,\sigma/\sqrt n\big)\approx0{,}95$.
\end{theoreme}

\begin{aretenir}
Une moyenne calculée sur $n$ observations se trompe typiquement de $\sigma/\sqrt n$ : pour diviser l'erreur par $2$, il faut $4$ fois plus de données.
C'est ce qui permet d'estimer l'erreur d'un modèle sur un échantillon test (chapitre~1).
\end{aretenir}

\subsection*{Exercices}

\begin{exercice}[Variable de Bernoulli]
\label{ex:bernoulli}
Soit $X\sim\mathcal B(p)$. Calculer $\E[X]$, $\E[X^2]$ et retrouver $\Var(X)=p(1-p)$. Pour quelle valeur de $p$ la variance est-elle maximale ?
\end{exercice}

\begin{exercice}[Précision d'un sondage]
\label{ex:sondage}
On estime la proportion $p$ de ménages équipés d'un réfrigérateur par la proportion observée sur $n$ ménages.
En utilisant le théorème central limite et $p(1-p)\le\frac14$, quelle taille $n$ garantit une erreur d'au plus $3$ points avec probabilité $95\,\%$ ?
\end{exercice}

\begin{exercice}[Meilleure constante]
\label{ex:meilleureconstante}
Un dé équilibré donne $Z\in\{1,\dots,6\}$. Calculer $\E[(Z-c)^2]$ pour $c=3$ et pour $c=3{,}5$, et vérifier la proposition~\ref{prop:meilleureconstante}.
\end{exercice}
